\documentclass[12pt, letterpaper]{article}

\usepackage{geometry}

\usepackage{float, graphicx}
    \graphicspath{ {image/} }

\usepackage{amsmath, amssymb, amsthm, cases, enumitem, nicefrac}
    \allowdisplaybreaks
    \theoremstyle{definition}
    \newtheorem{dfn}{Definition}
    \newtheorem{lem}{Lemma}
    \newtheorem{thm}{theorem}
    \newtheorem{cor}{corollary}
    \newtheorem{prp}{Proposition}

\usepackage{tikz}

\usepackage{enumitem}

\usepackage[toc]{multitoc}
    \setlength{\columnseprule}{0.5pt}

\usepackage{geometry}
    \geometry{letterpaper, margin = 1in}

\usepackage{hyperref}
    \hypersetup{
        colorlinks = true,
        allcolors = blue
    }


\begin{document}

\noindent
\begin{minipage}{0.5\textwidth}
    Will Daniels; 112774725

    April 22, 2023

    Computational Astrophysics; AST 390

    Homework 6
\end{minipage}
%
\begin{minipage}{0.5\textwidth}
    \begin{flushright}
        \includegraphics[height = 48pt]{../../LatexAssets/SBULogoStacked}
    \end{flushright}
\end{minipage}
\noindent
\rule{\textwidth}{1pt}

\section{Finite difference on Burgers' equation}

For this assigment, we consider two first-order finite-difference techniques for
solving Burgers' equation, \( u_t + u u_x = 0 \). The two analytically identical
forms are
\begin{equation}
  u_i^{n+1} = u_i^n -
  \frac{\Delta t}{\Delta x} u_i^n \left( u_i^n - u_{i-1}^n \right)
  \label{eq:method_a}
\end{equation}
and
\begin{equation}
  u_i^{n+1} = u_i^n + -
  \frac{\Delta t}{2 \Delta x} \left(
    \left( u_i^n \right)^2 - \left( u_{i-1}^n \right)^2
    \label{eq:method_b}
  \right)
\end{equation}
We will call \eqref{eq:method_a} method a and \eqref{eq:method_b} method b. The
numerical performance of the two methods is displayed in figure \ref{fig:plots}
and the shock speed tabulated in table \ref{tab:speed}. The simulation time is
0.1 seconds and the parameter \(C\) was choosen to be 0.5. It is true that for
32 and 64 bins the shock speed was exactly the same. This is because there are
not enough points to exactly capture where the shock is. This would improve with
more bins, but even 256 bins takes prohibitively long to run.

\begin{figure} [!b]
  \centering
  \caption{Plots at various number of points}
  \begin{subfigure}{0.3\textwidth}
    \centering
    \caption{32 points}
    \includegraphics[width=\textwidth]{../data/32.jpeg}
  \end{subfigure} \hfill \begin{subfigure}{0.3\textwidth}
    \centering
    \caption{64 points}
    \includegraphics[width=\textwidth]{../data/64.jpeg}
  \end{subfigure} \hfill \begin{subfigure}{0.3\textwidth}
    \centering
    \caption{128 points}
    \includegraphics[width=\textwidth]{../data/128.jpeg}
  \end{subfigure}
  \label{fig:plots}
\end{figure}

\begin{table} [!b]
  \centering
  \caption{Shock speeds}
  \begin{tabular}{c c c}
    \hline\hline
    Number of points & Method a & Method b \\
    \hline
    32  & 1.613 & 1.613 \\
    64  & 1.429 & 1.429 \\
    128 & 1.417 & 1.496 \\
    \hline\hline
  \end{tabular}
  \label{tab:speed}
\end{table}

We imediatly notice that method b predicts a faster shock speed than method a,
which raises the question which is correct, since the disritizations are
analytically equivalent. Although equivalent in the limit when the number of
discrete \(x\) values goes to infinity, method b has a stronger sense of
physicsal conservation laws, so we should trust it as being more accurate.

\end{document}
